% GENERATED FILE. Edit canonical lessons, then run npm run generate:books.
% canonical-source-hash: fnv1a64:3b5d030cdbdda014
% canonical-lessons: ZH-C535-yuanyi, ZH-W535-yuan-wish, ZH-C535-yuanwang, ZH-C535-ruyuan, ZH-C535-ziyuan, ZH-C535-danyuan

\chapter{If You're Willing}
\label{ch:zh-if-you-re-willing}
\hlchaptermodality{535}

\begin{chapteropening}
\textbf{By the end of this chapter:} \emph{I can say I'm willing, talk about wishes, getting what I hoped for and doing something of my own free will, and say let's hope.}

Complete the last lesson of chapter 535: I can say I'm willing, talk about wishes, getting what I hoped for and doing something of my own free will, and say let's hope.
\end{chapteropening}

\section[\texorpdfstring{\zh{愿意} (yuànyì)}{愿意 (yuànyì)}]{\zh{愿意} (yuànyì) — to be willing (to)}

\label{lesson:ZH-C535-yuanyi}



\begin{quote}
Before the new one: say the Mandarin for to give one's word, then the Mandarin for may all go as you wish.
\end{quote}

\subsection*{You'll want to know: \zh{愿意}}
\textbf{\zh{愿意}} — \emph{yuànyì} — \textquotedblleft{}to be willing (to)\textquotedblright{}.

\begin{quote}
\textbf{\zh{要是你愿意我们明天就去}} \emph{yàoshi nǐ yuànyì wǒmen míngtiān jiù qù} — If you're willing, we'll go tomorrow.
\end{quote}

The new character, in pieces: \zh{愿} — fourteen strokes in four pieces: a bar and a long stroke falling to the left, like a cliff; under it \zh{白}, five strokes, and \zh{小}, three strokes; and \zh{心} at the foot, four strokes. The next lesson puts it in your hand.

\begin{cousinweb}
\textbf{\zh{愿意}} (\emph{yuànyì}) is built from \zh{愿} and \zh{意} 'what one means', both inherited from Middle Chinese, \emph{ngjwonH} and \emph{'iH} (Guangyun, in Baxter's transcription), which glosses the first 'to desire' and 'to think of'.
\end{cousinweb}

\begin{grammarlens}[title={What you've built}]
One more word for saying what would happen.
\end{grammarlens}

\subsection*{Your turn}
\begin{itemize}
\raggedright
  \item \emph{Say it:} \emph{yuànyì}
  \item \emph{Say it:} \emph{yuànyì}, once more
  \item \emph{Say it:} \emph{yàoshi nǐ yuànyì wǒmen míngtiān jiù qù} — If you're willing, we'll go tomorrow.
  \item \emph{Recall:} say the Mandarin for to give one's word, then the Mandarin for may all go as you wish, then say \emph{yuànyì} again
\end{itemize}

\begin{tcolorbox}[breakable,colback=teal!4,colframe=teal!35!black,title={Before you move on}]
What does \zh{愿意} mean? (\textquotedblleft{}To be willing (to)\textquotedblright{}.) Say it once more.
\end{tcolorbox}
\section[\texorpdfstring{\zh{愿} (yuàn)}{愿 (yuàn)}]{\zh{愿} — a cliff over \zh{白} and \zh{小}, then \zh{心}}

\label{lesson:ZH-W535-yuan-wish}



\begin{quote}
The word first, then the character inside it.

\begin{itemize}
\raggedright
  \item \emph{Recall:} say \emph{yuànyì} — \textbf{R1}, one lesson back, before it decays
\end{itemize}

\begin{quote}
\zh{白} \zh{心}
\end{quote}

\emph{Read it:} each one aloud, with nothing said first — \emph{bái}, \emph{xīn}
\end{quote}

\subsection*{Script — a cliff over \zh{白} and \zh{小}, then \zh{心}}
\hlblockfigure{\detokenize{figures/ZH-W535-yuan-wish-filmstrip.pdf}}{How \zh{愿} is written, stroke by stroke}

\begin{quote}
\zh{愿}
\end{quote}

\textbf{Four pieces: the cliff, then what it shelters, then \zh{心} at the foot.} Count the pieces, not the strokes.

\textbf{The cliff, two strokes}: a bar, then a long stroke falling to the left down from its left end.

\textbf{Under the cliff, \zh{白}, five strokes, then \zh{小}, three strokes}, as you wrote them: a short stroke falling to the left, a vertical, a bar that turns down, two short bars; then a vertical that hooks to the left, a short stroke on its left and a dot on its right.

\textbf{\zh{心}, four strokes}, as you wrote it, across the foot: a dot, a long curve that hooks up, then two dots.

2 + 5 + 3 + 4 = fourteen.

\subsection*{Your turn}
\begin{itemize}
\raggedright
  \item \emph{Write it:} \zh{愿}, all fourteen, with the model in front of you
  \item \emph{Cover:} write it again from memory, and count the pieces, not the strokes
\end{itemize}

\begin{tcolorbox}[breakable,colback=teal!4,colframe=teal!35!black,title={Before you move on}]
How many pieces does \zh{愿} have? (\textbf{Four: the cliff, \zh{白}, \zh{小} and \zh{心}}.) Which piece comes last? (\textbf{\zh{心}, at the foot}.)

Source: \href{https://raw.githubusercontent.com/chanind/hanzi-writer-data/68d10a4b21150cae5e1ebbd223eed289cf32d90c/data/\%E6\%84\%BF.json}{Hanzi Writer Data \zh{愿}, pinned stroke paths and medians}.
\end{tcolorbox}
\section[\texorpdfstring{\zh{愿望} (yuànwàng)}{愿望 (yuànwàng)}]{\zh{愿望} (yuànwàng) — a wish one has; a desire}

\label{lesson:ZH-C535-yuanwang}



\begin{quote}
Before the new one: say the Mandarin for may all go as you wish, then the Mandarin for to be willing.
\end{quote}

\subsection*{You'll want to know: \zh{愿望}}
\textbf{\zh{愿望}} — \emph{yuànwàng} — \textquotedblleft{}a wish one has; a desire\textquotedblright{}.

\begin{quote}
\textbf{\zh{如果可以有一个愿望你想要什么}} \emph{rúguǒ kěyǐ yǒu yí ge yuànwàng nǐ xiǎng yào shénme} — If you could have one wish, what would you want?
\end{quote}

\begin{cousinweb}
\textbf{\zh{愿望}} (\emph{yuànwàng}) is built from \zh{愿} 'to desire' and \zh{望} 'a hope', both inherited from Middle Chinese, \emph{ngjwonH} and \emph{mjangH} (Guangyun, in Baxter's transcription).
\end{cousinweb}

\begin{grammarlens}[title={What you've built}]
One more word for saying what would happen.
\end{grammarlens}

\subsection*{Your turn}
\begin{itemize}
\raggedright
  \item \emph{Say it:} \emph{yuànwàng}
  \item \emph{Say it:} \emph{yuànwàng}, once more
  \item \emph{Say it:} \emph{rúguǒ kěyǐ yǒu yí ge yuànwàng nǐ xiǎng yào shénme} — If you could have one wish, what would you want?
  \item \emph{Recall:} say the Mandarin for may all go as you wish, then the Mandarin for to be willing, then say \emph{yuànwàng} again
\end{itemize}

\begin{tcolorbox}[breakable,colback=teal!4,colframe=teal!35!black,title={Before you move on}]
What does \zh{愿望} mean? (\textquotedblleft{}A wish one has; a desire\textquotedblright{}.) Say it once more.
\end{tcolorbox}
\section[\texorpdfstring{\zh{如愿} (rúyuàn)}{如愿 (rúyuàn)}]{\zh{如愿} (rúyuàn) — to get what one wished for}

\label{lesson:ZH-C535-ruyuan}



\begin{quote}
Before the new one: say the Mandarin for to be willing, then the Mandarin for a wish one has.
\end{quote}

\subsection*{You'll want to know: \zh{如愿}}
\textbf{\zh{如愿}} — \emph{rúyuàn} — \textquotedblleft{}to get what one wished for\textquotedblright{}.

\begin{quote}
\emph{zhǐyào nǐ nǔlì jiù néng rúyuàn} — As long as you work hard, you'll get what you hoped for.
\end{quote}

\begin{cousinweb}
\textbf{\zh{如愿}} (\emph{rúyuàn}) is built from \zh{如} 'as, like' and \zh{愿} 'a wish', both inherited from Middle Chinese, \emph{nyo} and \emph{ngjwonH} (Guangyun, in Baxter's transcription): as one wished.
\end{cousinweb}

\begin{grammarlens}[title={What you've built}]
One more word for saying what would happen.
\end{grammarlens}

\subsection*{Your turn}
\begin{itemize}
\raggedright
  \item \emph{Say it:} \emph{rúyuàn}
  \item \emph{Say it:} \emph{rúyuàn}, once more
  \item \emph{Say it:} \emph{zhǐyào nǐ nǔlì jiù néng rúyuàn} — As long as you work hard, you'll get what you hoped for.
  \item \emph{Recall:} say the Mandarin for to be willing, then the Mandarin for a wish one has, then say \emph{rúyuàn} again
\end{itemize}

\begin{tcolorbox}[breakable,colback=teal!4,colframe=teal!35!black,title={Before you move on}]
What does \zh{如愿} mean? (\textquotedblleft{}To get what one wished for\textquotedblright{}.) Say it once more.
\end{tcolorbox}
\section[\texorpdfstring{\zh{自愿} (zìyuàn)}{自愿 (zìyuàn)}]{\zh{自愿} (zìyuàn) — of one's own free will; voluntary}

\label{lesson:ZH-C535-ziyuan}



\begin{quote}
Before the new one: say the Mandarin for a wish one has, then the Mandarin for to get what one wished for.
\end{quote}

\subsection*{You'll want to know: \zh{自愿}}
\textbf{\zh{自愿}} — \emph{zìyuàn} — \textquotedblleft{}of one's own free will; voluntary\textquotedblright{}.

\begin{quote}
\textbf{\zh{我们自愿在这儿工作}} \emph{wǒmen zìyuàn zài zhèr gōngzuò} — We work here of our own free will.
\end{quote}

\begin{cousinweb}
\textbf{\zh{自愿}} (\emph{zìyuàn}) is built from \zh{自} 'self' and \zh{愿} 'to wish', both inherited from Middle Chinese, \emph{dzijH} and \emph{ngjwonH} (Guangyun, in Baxter's transcription), which glosses the first 'from' and 'by way of'.
\end{cousinweb}

\begin{grammarlens}[title={What you've built}]
One more word for saying what would happen.
\end{grammarlens}

\subsection*{Your turn}
\begin{itemize}
\raggedright
  \item \emph{Say it:} \emph{zìyuàn}
  \item \emph{Say it:} \emph{zìyuàn}, once more
  \item \emph{Say it:} \emph{wǒmen zìyuàn zài zhèr gōngzuò} — We work here of our own free will.
  \item \emph{Recall:} say the Mandarin for a wish one has, then the Mandarin for to get what one wished for, then say \emph{zìyuàn} again
\end{itemize}

\begin{tcolorbox}[breakable,colback=teal!4,colframe=teal!35!black,title={Before you move on}]
What does \zh{自愿} mean? (\textquotedblleft{}Of one's own free will; voluntary\textquotedblright{}.) Say it once more.
\end{tcolorbox}
\section[\texorpdfstring{\zh{但愿} (dànyuàn)}{但愿 (dànyuàn)}]{\zh{但愿} (dànyuàn) — let us hope that; if only}

\label{lesson:ZH-C535-danyuan}



\begin{quote}
Before the new one: say the Mandarin for to get what one wished for, then the Mandarin for of one's own free will.
\end{quote}

\subsection*{You'll want to know: \zh{但愿}}
\textbf{\zh{但愿}} — \emph{dànyuàn} — \textquotedblleft{}let us hope that; if only\textquotedblright{}.

\begin{quote}
\textbf{\zh{但愿他明天能来}} \emph{dànyuàn tā míngtiān néng lái} — Let's hope he can come tomorrow.
\end{quote}

\begin{cousinweb}
\textbf{\zh{但愿}} (\emph{dànyuàn}) is built from \zh{但} 'only' and \zh{愿} 'to wish', both inherited from Middle Chinese, \emph{danX} and \emph{ngjwonH} (Guangyun, in Baxter's transcription), which glosses the first as a grammar word and 'merely': I only wish.
\end{cousinweb}

\begin{grammarlens}[title={What you've built}]
One more word for saying what would happen.
\end{grammarlens}

\subsection*{Your turn}
\begin{itemize}
\raggedright
  \item \emph{Say it:} \emph{dànyuàn}
  \item \emph{Say it:} \emph{dànyuàn}, once more
  \item \emph{Say it:} \emph{dànyuàn tā míngtiān néng lái} — Let's hope he can come tomorrow.
  \item \emph{Recall:} say the Mandarin for to get what one wished for, then the Mandarin for of one's own free will, then say \emph{dànyuàn} again
\end{itemize}

\begin{tcolorbox}[breakable,colback=teal!4,colframe=teal!35!black,title={Before you move on}]
What does \zh{但愿} mean? (\textquotedblleft{}Let us hope that; if only\textquotedblright{}.) Say it once more.
\end{tcolorbox}
